% Integrated 4 September 2026; exact in-project provenance is recorded below.
% The exact fixed level, lattice, chamber, height, theorem hypotheses, and locators are retained.
\section{Primary-source theorem audit of effective counting for the fixed Koide ternary quadric sector}\label{koide-lit-primary-source-theorem-audit-of-effective-counting-for-the-fixed-koide-ternary-quadric-sector}

\textbf{Date:} 4 September 2026\\
\textbf{Exact counting function:}

\[
 R(X)=\#\left\{(a,q,n)\in\mathbb Z^3:
 \begin{array}{l}
 a^2-3q^2-18n^2=-2,\\
 a>0,\ q\geq0,\ n\geq0,\ a+4n>q,\\
 4a+17n\leq X
 \end{array}\right\}.
\]

The leading constant proved in
\cref{koide-count-thm-linear-asymptotic} is

\[
 C_0=\frac{9}{2\pi^2}\log(2+\sqrt3)
       \log\!\left(\frac{3}{2\sqrt2}\right).
\]

The canonical project-relative mathematical sources used by this audit are:

\begin{itemize}
\item
  \path{supporting_materials/workbench/research/koide_counting_repair.tex},
  which proves the coordinate bijection, local densities, invariant-volume
  factor, Gorodnik--Oh--Shah hypotheses, and calibrated constant, SHA-256
  \path{965b0a3dca8bfe4f22b9e9208368fa04697070fd25bd06bd67f2481d17cb9d9d};
\item
  \path{supporting_materials/workbench/research/koide_exact_parametrization.tex},
  which proves the exact parametrization and the Benoist--Oh power-saving
  application, SHA-256
  \path{a457709dcce0fce0774eb9326b84b00d5d8144c2087f480da704974e499bcdd3};
\item
  \path{supporting_materials/workbench/research/koide_quaternion_bridge.tex},
  which proves the exact order, shell, centralizer, conjugacy-class, and
  radial-height statements used below, SHA-256
  \path{93759f7ba99870c84c9144a373a53d2e9a8bc3ae909518e47991373ace0b47bb};
\item
  \path{supporting_materials/workbench/research/koide_exact_census.tex},
  which records the separately bounded finite census, SHA-256
  \path{18ed5968b7e67d9a0b23d8675b40c37ba46170ad29efdbb9ce10ed632c99600a}.
\end{itemize}

The primary publications themselves are identified by exact versions,
dates where version-sensitive, locators, and DOI or arXiv links in
\cref{koide-lit-primary-references-and-exact-locators}.  They are not
claimed to be stored in a separate canonical evidence manifest.

\subsection{Determination}\label{koide-lit-determination}

Among the source versions enumerated in
\cref{koide-lit-primary-references-and-exact-locators}, no theorem proves,
for this exact fixed level, full integral lattice, sharp chamber, and sharp
polyhedral height, the unconditional estimate

\[
 R(X)-C_0X=O_\varepsilon(X^{1/2+\varepsilon}).
\]

The strongest exact-domain conclusion obtained from the general theorems
enumerated there is

\[
 \boxed{R(X)=C_0X+O(X^{1-\eta})\quad\text{for some unspecified }\eta>0.}
\]

It follows by applying Benoist--Oh, Definition 13.1 and Theorem 13.2, orbit by orbit after directly verifying effective well-roundedness for the present family. Benoist--Oh do not give a numerical lower bound for \(\eta\), and their proof does not imply \(\eta\geq 1/2\).

The superficially stronger exponents in nearby problems do not answer the question:

\begin{itemize}

\item
  Good's sharp hyperbolic conjugacy-class theorem, as formulated and reproved by Chatzakos--Petridis, gives \(O(X^{2/3})\) only for a different radial cutoff and only after every small-eigenvalue term \(X^{s_j}\) has been placed in the main expression. The same paper states \(O_\varepsilon(X^{1/2+\varepsilon})\) as Conjecture 5.7.
\item
  Liu--Sarnak obtain a smooth radial error \(O_\varepsilon(X^{1/2+\theta+\varepsilon})\), hence \(O_\varepsilon(X^{39/64+\varepsilon})\) from \(\theta=7/64\), and square-root strength only under the generalized Selberg/Ramanujan-at-infinity conjecture for the relevant Jacquet--Langlands transfers. Their paper's global affine-sieve setup assumes \(t\det(f)\) square-free, but the analytic Proposition 3.2 itself is formulated for a variable cocompact lattice with a spectral gap. After the exact quaternionic identification it applies class by class to this lattice. What it still does not supply is a sharp cutoff, the coordinate chamber, or the nonradial height \(4a+17n\).
\item
  Gorodnik--Nevo's published sharp polar-sector theorem loses powers through set thickening. For \(\dim G=3\), it gives \(O(X^{19/20+\epsilon})\) even with tempered spectrum; the Kim--Sarnak complementary-series threshold, which does apply to the exact congruence arithmetic lattice by Jacquet--Langlands, gives the still weaker \(O(X^{39/40+\epsilon})\). This remains a comparison theorem because Gorodnik--Nevo's fixed polar-sector family has not been identified with the exact polyhedron, including all integral orbits.
\item
  Square-root-scale statements in Chatzakos--Petridis, Chatzakos, Voskou, and Lekkas--Petridis are mean-square, averaged, Omega, or conjectural statements. They are not pointwise sharp-sector upper bounds.
\item
  Recent split-stabilizer theorems assume \(-m\det(F)\) is a square and have a \(T\log T\)-type main term. Here \(-(-2)\det(F)=108\) is not a square, so they concern the opposite arithmetic regime.
\end{itemize}

The two independent missing inputs are therefore (i) a sharp-boundary spectral or arithmetic estimate for the actual height and chamber, and (ii) control or elimination of any complementary-spectrum secondary terms for the actual sector. The Pell faces \(q=0\) and \(n=0\), which contribute only \(O(\log X)\), do not solve the moving outer-boundary problem.

\subsection{Exact arithmetic and geometric classification}\label{koide-lit-exact-arithmetic-and-geometric-classification}

\subsubsection{The rational group and stabilizer}\label{koide-lit-the-rational-group-and-stabilizer}

Let

\[
 F(a,q,n)=a^2-3q^2-18n^2,
 \qquad \mathscr X=\{F=-2\}.
\]

The matrix of \(F\) is \(\operatorname{diag}(1,-3,-18)\), so \(\det(F)=54\) and

\[
 -(-2)\det(F)=108\notin\mathbb Q^{\times2}.
\]

The homogeneous form \(F\) is anisotropic over \(\mathbb Q\). Indeed, suppose a primitive integral triple satisfies

\[
 a^2-3q^2-18n^2=0.
\]

Modulo \(3\), one gets \(3\mid a\); write \(a=3a_1\). Dividing by \(3\) and reducing modulo \(3\) gives \(3\mid q\); write \(q=3q_1\). The resulting equation is

\[
 a_1^2=3q_1^2+2n^2.
\]

Modulo \(3\), this says \(a_1^2+n^2\equiv0\pmod3\), whence \(3\mid a_1,n\). Thus \(3\mid a,q,n\), contradicting primitivity.

Put \(G=\operatorname{Spin}(F)\) and \(x_0=(1,1,0)\). The orthogonal complement of \(x_0\) is

\[
 x_0^\perp=\{(3s,s,t):s,t\in\mathbb Q\},
 \qquad F(3s,s,t)=6(s^2-3t^2).
\]

Consequently the connected stabilizer \(H=\operatorname{Stab}_G(x_0)^\circ\) is the norm-one torus

\[
 \operatorname{Res}^{1}_{\mathbb Q(\sqrt3)/\mathbb Q}\mathbb G_m.
\]

It is split over \(\mathbb R\), anisotropic over \(\mathbb Q\), and has no
nontrivial \(\mathbb Q\)-character.  Under the quaternionic identification
in \cref{koide-quat-centralizer-optimal-quadratic-order-and-primitivity},
\(x_0\) corresponds to \(v_0=(I+J)/2\), with \(v_0^2=3\), and

\[
 \mathcal O\cap\mathbb Q(v_0)=\mathbb Z[v_0],
 \qquad
 H(\mathbb Q)\cap\Gamma
 =\{\pm(2+v_0)^k:k\in\mathbb Z\}.
\]

Thus \((H(\mathbb Q)\cap\Gamma)/\{\pm1\}\) is infinite cyclic.  Under
\(H(\mathbb R)^\circ\cong\mathbb R_{>0}\), the logarithm sends its positive
part to \(\log(2+\sqrt3)\mathbb Z\subset\mathbb R\); the quotient is compact,
and in particular \((H\cap\Gamma)\backslash H(\mathbb R)\) has finite
volume. Since \(F\) is \(\mathbb Q\)-anisotropic, the arithmetic Fuchsian
quotient attached to \(G\) is cocompact as well. There is no ambient
hyperbolic cusp.

This is the finite-volume stabilizer regime associated with a closed geodesic.
It is not the rationally split regime with divergent stabilizer in which one
obtains \(T\log T\) growth.

\subsubsection{The exact height really is a norm-sector height}\label{koide-lit-the-exact-height-really-is-a-norm-sector-height}

Define

\[
 \|(a,q,n)\|_*=max\{|4a+17n|,|q|,|a|\}.
\]

This is a genuine norm because it is the pullback of the sup norm under the invertible linear map

\[
 (a,q,n)\longmapsto(4a+17n,q,a),
\]

whose determinant is \(-17\). On the target chamber,

\[
 0<a<4a+17n,
 \qquad 0\leq q<a+4n<4a+17n,
\]

and therefore \(\|(a,q,n)\|_*=4a+17n\). Thus the outer cutoff is not merely comparable to a norm: it is exactly the restriction of a fixed ambient norm.

The additional wall \(a+4n>q\) removes only a bounded real set. If \(q\geq a+4n\), then

\[
 a^2+2=3q^2+18n^2
 \geq3(a+4n)^2+18n^2,
\]

so \(2a^2+24an+66n^2\leq2\). Its only integral point in the positive orthant on \(\mathscr X\) is \((1,1,0)\). The unbounded faces \(q=0\) and \(n=0\) are Pell families and contribute \(O(\log X)\) integral points, but their real boundary strips still have to be included in a sharp-domain argument.

\subsubsection{Invariant volume with a strong enough secondary estimate}\label{koide-lit-invariant-volume-with-a-strong-enough-secondary-estimate}

Use

\[
 \sqrt3q=r\cos\theta,
 \qquad 3\sqrt2n=r\sin\theta,
 \qquad 0\leq\theta\leq\frac\pi2.
\]

On \(F=-2\), \(r=\sqrt{a^2+2}\), and the gauge measure is

\[
 dm_\infty=\frac1{6\sqrt6}\,da\,d\theta.
\]

Writing \(c=17/(3\sqrt2)\), the outer boundary is

\[
 4a+c\sqrt{a^2+2}\sin\theta=X.
\]

If \(A_X(\theta)\) is its positive solution, then uniformly in \(\theta\),

\[
 A_X(\theta)=\frac{X}{4+c\sin\theta}+O(X^{-1}).
\]

Indeed \(A_X(\theta)\asymp X\) uniformly and \(\sqrt{A^2+2}=A+O(A^{-1})\). Once \(X\) contains the bounded wall excision, it follows that

\[
 m_\infty(B_X)=I_\infty X+O(1),
 \qquad
 I_\infty=\frac1{2\sqrt3}\log\!\left(\frac3{2\sqrt2}\right).
\]

This \(O(1)\) refinement matters: a theorem whose main term is the exact invariant volume can be converted to the displayed \(C_0X\) main term without spending a power of \(X\).

\subsubsection{The full integral locus is not assumed to be one orbit}\label{koide-lit-the-full-integral-locus-is-not-assumed-to-be-one-orbit}

The exact lattice must not be silently replaced by a convenient single orbit or by a commensurable maximal-order lattice. Borel--Harish-Chandra finiteness, in the form quoted by Benoist--Oh on journal p.~1892, says that \(G(\mathbb Z)\) has finitely many orbits on \(\mathscr X(\mathbb Z)\). Thus

\[
 \mathscr X(\mathbb Z)=\bigsqcup_{j=1}^{h}x_j\Gamma
\]

for the integral spin group \(\Gamma\) attached to the chosen lattice and finitely many integral representatives \(x_j\). An effective theorem for one orbit must be applied to every \(x_j\Gamma\), and the finitely many main-measure scalars must be summed. The full-lattice scalar is calibrated from Borovoi's local/Brauer--Manin density formula in \cref{koide-count-prop-finite-sign-balance,koide-count-thm-linear-asymptotic}; it is not obtained by declaring that all integral points form the orbit of \(x_0\).

\subsubsection{Exact quaternionic lattice and its congruence status}\label{koide-lit-exact-quaternionic-lattice-and-its-congruence-status}

For the spectral comparisons below, the precise arithmetic lattice can be retained. Put

\[
 B=(18,-6)_{\mathbb Q}
   =\mathbb Q\langle I,J\rangle/(I^2-18,\ J^2+6,\ IJ+JI),
 \qquad K=IJ,
\]

and

\[
 \mathcal O=\mathbb Z1\oplus
  \mathbb Z\frac{I+J}{2}\oplus
  \mathbb Z\frac{I-J}{2}\oplus
  \mathbb Z\frac K2.
\]

The complete integral multiplication table is
\eqref{koide-quat-eq-multiplication-table}, and its reduced-trace Gram
determinant is \(-108^2\) by
\eqref{koide-quat-eq-order-discriminant-determinant}; hence
\(\operatorname{discrd}(\mathcal O)=108\).  The Hilbert-symbol calculations
\eqref{koide-quat-eq-hilbert-three} and
\eqref{koide-quat-eq-hilbert-two}, together with the unramified-place and
real-place calculation summarized in
\eqref{koide-quat-eq-ramification-discriminant}, show that the ramification
set is exactly \(\{2,3\}\).  Thus \(B\) is the indefinite quaternion
division algebra of discriminant \(6\). The map

\[
 (a,q,n)\longmapsto
 \gamma(a,q,n)=2+\frac{qI+aJ+nK}{2}
\]

is the bijection \eqref{koide-quat-eq-shell-bijection} from
\(F(a,q,n)=-2\) to the reduced-trace-\(4\), reduced-norm-one locus in
\(\mathcal O\). Indeed, the parity condition \(a\equiv q\pmod2\), which is
exactly what membership in \(\mathcal O\) requires, follows automatically
from \(F=-2\) modulo \(2\), and

\[
 \operatorname{nrd}(\gamma)=4+\frac32F(a,q,n)=1.
\]

Thus the acting Fuchsian lattice is exactly \(\Gamma=\mathcal O^1\), not a substituted maximal-order lattice. If \(\mathcal M\supset\mathcal O\) is maximal, the reduced-discriminant index formula gives

\[
 [\mathcal M:\mathcal O]=\frac{108}{6}=18,
 \qquad 18\mathcal M\subset\mathcal O,
\]

and consequently

\[
 \mathcal M^1(18\mathcal M)
 \subset\mathcal O^1\subset\mathcal M^1.
\]

Here
\(\mathcal M^1(18\mathcal M)=\{u\in\mathcal M^1:u-1\in18\mathcal M\}\).
Therefore \(\mathcal O^1\) is a congruence arithmetic Fuchsian group in the
exact sense used by Chu--Li. The number \(108\) is the order's reduced
discriminant, not an Eichler level: the order is nonmaximal at the ramified
primes \(2,3\). This non-Eichler feature changes the finite compact-open
subgroup and hence which automorphic representations have fixed vectors and
the dimensions of those fixed-vector spaces; it does not change the
Jacquet--Langlands/Kim--Sarnak archimedean bound.

\subsection{Result matrix}\label{koide-lit-result-matrix}

\begin{longtable}[]{@{}
  >{\raggedright\arraybackslash}p{0.31\linewidth}
  >{\raggedright\arraybackslash}p{0.31\linewidth}
  >{\raggedright\arraybackslash}p{0.31\linewidth}@{}}
\toprule\noalign{}
\begin{minipage}[b]{\linewidth}\raggedright
Source and exact locator
\end{minipage} & \begin{minipage}[b]{\linewidth}\raggedright
What is actually proved
\end{minipage} & \begin{minipage}[b]{\linewidth}\raggedright
Status for the target
\end{minipage} \\
\midrule\noalign{}
\endhead
\bottomrule\noalign{}
\endlastfoot
Duke--Rudnick--Sarnak (1993), Theorem 1.2, p.~145; Example 1.5, pp.~146--147 & One arithmetic orbit on an affine symmetric variety: count \(\sim\) invariant volume; ternary anisotropic stabilizer gives linear rather than \(T\log T\) growth & Correct qualitative regime, no rate and no exact sharp-sector theorem \\
Eskin--McMullen (1993), Theorem 1.4, author PDF p.~3 & Every well-rounded sequence has count asymptotic to volume & Qualitative only; no exponent \\
Gorodnik--Oh--Shah (2009), definitions before Theorem 1.6 and Theorem 1.6, p.~5; Theorem 1.18, pp.~8--9 & Any norm and an admissible generic cone: full integral count \(\sim\) volume & Direct qualitative theorem for the exact norm and cone, whose hypotheses are proved in \cref{koide-count-lem-genuine-norm,koide-count-lem-gos-hypotheses}; no error term \\
Gorodnik--Oh--Shah (2010), Theorem 1.13, author PDF pp.~4--5 & Well-rounded polar sectors and their asymptotic & Qualitative only \\
Benoist--Oh (2012), Theorem 1.3, p.~1894 & Standard product/Euclidean height ball: relative \(O(T^{-\delta})\) for some \(\delta>0\) & Theorem 1.3 alone is not the exact polyhedron \\
Benoist--Oh (2012), Definition 13.1, p.~1931, and Theorem 13.2, pp.~1931--1932 & Any independently verified effectively well-rounded family: \(\operatorname{vol}(B)(1+O(\operatorname{vol}(B)^{-\delta_0}))\) & Applies to the exact family after the proof in \cref{koide-lit-definition-13.1-for-the-actual-b_x}; yields only unspecified \(O(X^{1-\eta})\) \\
Gorodnik--Nevo (2012), Definition 3.1, p.~13; Theorem 8.1 and Corollary 8.3, pp.~34--35 & Sharp fixed polar sectors: \(O(\operatorname{vol}^{1-\zeta+\epsilon})\), \(\zeta=(2n_e(p))^{-1}(1+3\dim G)^{-1}\) & Explicit but much weaker; target polyhedron and finite orbit union require an additional reduction \\
Liu--Sarnak (2010), equations (2.3)--(2.8), pp.~4--5; Proposition 3.2, pp.~10--11; equation (3.65), p.~16 & Smooth radial count for a cocompact lattice containing the fixed cyclic stabilizer: \(O_\epsilon(T^{1/2+\theta+\epsilon})\) & Exact lattice and finite class sum are available; unconditional \(39/64\), conditional \(1/2\), but fixed-width smoothing and the radial height do not match the sharp target \\
Good (1983), Theorem 4, p.~116, as explicitly restated and reproved in Chatzakos--Petridis (2016), Theorem 1.1, p.~3 & Sharp radial hyperbolic conjugacy-class count: \(O(X^{2/3})\) after all \(X^{s_j}\) small-spectrum terms & Closest sharp spectral model, not the exact polyhedron; square root is CP Conjecture 5.7 \\
Bourgain--Kontorovich--Sarnak (2010), Theorem 1.1, p.~2 & Fixed angular Fourier modes in a sharp group norm ball: \(O_{n,k}(T^{4/3})\) after \(T^{2s_j}\) terms & With affine height \(X\asymp T^2\), again \(X^{2/3}\); constants are not uniform enough to recover a sharp interval for free \\
Browning--Gorodnik (2017), Theorem 2.1, p.~11; Corollaries 2.5--2.6, pp.~24--25 & Full Euclidean balls, finite adelic/orbit sums, and congruence conditions: unspecified \(O(\operatorname{vol}^{1-\varrho})\) & Correct nonsquare homogeneous regime, wrong height geometry and no explicit exponent \\
Cao--Huang (2022), Theorem 1.5, p.~4; Theorem 3.4, pp.~21--22 & Arithmetic purity for anisotropic nonsquare ternary quadrics; effective input for Euclidean balls has an unspecified power & Qualitative full-lattice theorem, not a square-root sector result \\
Cao--Huang--Zhang (arXiv:2509.23320), Setting 2.1 and Theorem 2.1 & Product Euclidean/max height balls with congruence-uniform \(O(\mathbf N(\mathfrak l)^\kappa\operatorname{vol}^{1-\delta})\) & Current extension, but \(\delta\) is unspecified and the growing real factor is the standard height ball, not this cone \\
Oh--Shah (2014), Theorem 1.2, p.~3; Theorem 6.1, pp.~13--14 & Divergent-geodesic/split-stabilizer \(T\log T\) count; smooth error \(O(T)\), sharp relative logarithmic error & Excluded by the finite-volume anisotropic stabilizer \\
Xu--Zhang (2023), main theorem; Huang (arXiv:2412.03350), Theorem 1.1; Zhang (arXiv:2603.21080), Theorem 1.1 & Split-stabilizer results assuming \(-m\det(F)\) is a square, with smooth or radial \(T\log T\)-type expansions & \(108\) is nonsquare; wrong main-term regime \\
Huang (arXiv:2412.05903), Theorems 1.1--1.2 & Smooth compact real weight while the level varies as \(m_0p_0^{2h}\); error of logarithmic relative strength & Varying level and smooth weight, not the fixed-level sharp count \\
Li--Liu (arXiv:2501.15033v2), Theorems 1.1--1.2, pp.~2--3; Lemma 2.3, pp.~5--6 & Almost-prime results using a smooth radial orbit count with \(O_\epsilon(d^{1+\epsilon}T^{1/2+\theta+\epsilon})\) & Current use of the Liu--Sarnak input; globally assumes \(t\det(f)\) square-free and supplies no sharp-sector count \\
Voskou (arXiv:2407.03134v3), Theorem 3 and Conjecture 1, p.~3 & Four sharp orientation sectors in a closed-geodesic double-coset problem: \(O(X^{2/3})\); square root conjectured & Current sharp comparator, not the target \\
Lekkas--Petridis (arXiv:2509.12902, version dated 24 Aug.~2026), Theorem 1.2, Theorem 1.5, Conjecture 1.6 & Sharp closed-geodesic double-coset \(O(X^{2/3})\), square-root mean square, pointwise square-root conjecture & Version-specific confirmation that mean square does not furnish the requested pointwise bound \\
\end{longtable}

\subsection{Benoist--Oh: exact applicability and exact limitation}\label{koide-lit-benoistoh-exact-applicability-and-exact-limitation}

Among the general effective theorems enumerated in
\cref{koide-lit-result-matrix}, this is the one whose stated hypotheses can
be verified for the actual family without changing its height or chamber.

\subsubsection{What the theorem assumes}\label{koide-lit-what-the-theorem-assumes}

Benoist--Oh work with a connected semisimple simply connected almost \(K\)-simple group \(G\), a finite-kernel rational representation, a symmetric subgroup \(H\), and the hypothesis that \(H^\circ\) has no nontrivial \(K\)-character; see journal pp.~1892--1893. For a congruence subgroup \(\Gamma_S\), Theorem 12.5, journal p.~1929, supplies the effective equidistribution of the translates of \((\Gamma_S\cap H_S)\backslash H_S\) used in Theorem 13.2. Example (A), journal p.~1896, explicitly treats ternary quadratic forms provided the homogeneous form does not represent zero over \(\mathbb Q\).

All of these hypotheses hold here:

\begin{itemize}

\item
  \(G=\operatorname{Spin}(F)\) is simply connected and absolutely almost simple of type \(A_1\), and \(G(\mathbb R)\) is noncompact.
\item
  The standard three-dimensional orthogonal representation has finite kernel.
\item
  The stabilizer of a nonisotropic vector is symmetric.
\item
  The mod-\(3\) descent in
  \cref{koide-lit-the-rational-group-and-stabilizer} proves that \(F\) does
  not represent zero over \(\mathbb Q\).
\item
  \(H^\circ\) is the \(\mathbb Q\)-anisotropic norm-one torus of \(\mathbb Q(\sqrt3)\), so it has no nontrivial \(\mathbb Q\)-character and its arithmetic quotient is compact.
\item
  With \(S=\{\infty\}\), the full integral spin group is itself a congruence subgroup in Benoist--Oh's sense.
\end{itemize}

There is no omitted wavefront or properness condition: symmetry of \(H\) is exactly what invokes their effective wavefront argument in Lemma 12.3 and Theorem 12.5. Cocompactness makes escape in the automorphic quotient simpler, but does not remove Definition 13.1(3), which still has to be checked on the growing affine set.

Here ``congruence'' is an exact property of the chosen integral group, not an inference from mere commensurability. Under the quaternionic realization of \(G\), the finite-integrality conditions define a compact-open subgroup of \(G(\mathbb A_f)\); equivalently, after placing the corresponding order \(\mathcal O\) in a maximal order \(\mathcal O_{\max}\), some principal subgroup \(\{u\in\mathcal O_{\max}^1:u\equiv1\pmod{N\mathcal O_{\max}}\}\) lies in the exact integral spin group. This distinction matters for spectral-gap transfer: an arbitrary finite-index noncongruence subgroup need not inherit the Kim--Sarnak gap.

\subsubsection{\texorpdfstring{Definition 13.1 for the actual \(B_X\)}{Definition 13.1 for the actual B\_X}}\label{koide-lit-definition-13.1-for-the-actual-b_x}

Let \(B_X\subset\mathscr X(\mathbb R)\) be the exact chamber cut out by the five linear conditions and \(4a+17n\leq X\). In Benoist--Oh's notation, \(B_{X,\epsilon}^+=B_XB_\epsilon\) and \(B_{X,\epsilon}^-=\bigcap_{u\in B_\epsilon}B_Xu\).

Condition (1) is void because there is no finite-place factor.

For condition (2), if \(u\) lies in an \(\epsilon\)-ball about the identity, every fixed linear functional \(L\) defining the chamber or height satisfies

\[
 |L(xu)-L(x)|\ll\epsilon\|x\|_2.
\]

On the unbounded portion of this chamber, \(\|x\|_2\asymp 1+a\asymp1+(4a+17n)\). Hence a membership change under \(u\) is confined to:

\begin{enumerate}
\def\labelenumi{\arabic{enumi}.}

\item
  an angular strip of width \(O(\epsilon)\) at \(q=0\) or \(n=0\);
\item
  an \(O(\epsilon X)\)-thick radial shell at \(4a+17n=X\); or
\item
  an \(O(\epsilon)\)-neighborhood of one of the compact boundary pieces \(a=0\) or \(a+4n=q\).
\end{enumerate}

Since \(dm_\infty=(6\sqrt6)^{-1}da\,d\theta\), their total measure is \(O(\epsilon X)\). Thus

\[
 m_\infty(B_{X,\epsilon}^+\setminus B_{X,\epsilon}^-)
 \ll \epsilon\,m_\infty(B_X),
\]

so Definition 13.1(2) holds with \(\kappa=1\).

For condition (3), Benoist--Oh's Euclidean height is comparable to \(1+a\) throughout a fixed group-neighborhood enlargement of the chamber. For every \(k>0\), uniformly in \(0<\epsilon<1\),

\[
 \int_{B_{X,\epsilon}^+}H_\infty(x)^{-k}\,dm_\infty(x)
 \ll_k\int_0^{O(X)}(1+a)^{-k}\,da.
\]

The right side is \(O(X^{1-k})\) for \(0<k<1\), \(O(\log X)\) for \(k=1\), and \(O(1)\) for \(k>1\). As \(m_\infty(B_X)\asymp X\), this is \(O(m_\infty(B_X)^{1-\delta(k)})\) for some \(\delta(k)>0\) in every case. Definition 13.1(3) is therefore satisfied.

Proposition 14.2 of Benoist--Oh verifies only their standard product height balls. It should not be cited as though it automatically covers this polyhedron. The preceding direct verification is the necessary bridge to Theorem 13.2.

\subsubsection{Finite-orbit passage and main-term normalization}\label{koide-lit-finite-orbit-passage-and-main-term-normalization}

Benoist--Oh state on journal p.~1892 that \(G(\mathcal O_S)\) has finitely many orbits on \(\mathscr X(\mathcal O_S)\). Apply Theorem 13.2 to each exact integral orbit. If \(\eta_j>0\) is the exponent for the \(j\)-th orbit, take \(\eta=\min_j\eta_j>0\). The finite sum has the form

\[
 R(X)=\kappa_{\mathrm{orb}}\,m_\infty(B_X)+O(X^{1-\eta}).
\]

Here \(\kappa_{\mathrm{orb}}\) is the sum of the orbit-volume normalization
constants.  The comparison with Borovoi's Theorem 1.4.1 and Proposition
1.5.1 is carried out in
\cref{koide-count-thm-linear-asymptotic}; specifically,
\eqref{koide-count-eq-singular-series},
\eqref{koide-count-eq-archimedean-asymptotic}, and
\eqref{koide-count-eq-final-constant} give

\[
 \kappa_{\mathrm{orb}}=\mathfrak S(\mathscr X)
 =\frac{9\sqrt3}{\pi^2}\log(2+\sqrt3),
 \qquad
 \kappa_{\mathrm{orb}} I_\infty=C_0.
\]

Together with \(m_\infty(B_X)=I_\infty X+O(1)\), this yields

\[
 R(X)=C_0X+O(X^{1-\eta}).
\]

This counts the actual integral lattice. No assertion that \(\mathscr X(\mathbb Z)\) is one orbit, and no replacement by a maximal order, is involved.

\subsubsection{Why the exponent remains unspecified}\label{koide-lit-why-the-exponent-remains-unspecified}

The proof of Benoist--Oh Theorem 13.2 balances an effective-equidistribution/Sobolev error and a boundary-thickening error:

\[
 c_1\epsilon^{-p}v_X^{1-\delta}+c_2\epsilon^\kappa v_X,
 \qquad v_X=\operatorname{vol}(B_X).
\]

The resulting saving is of the form

\[
 \eta=\frac{\delta\kappa}{p+\kappa}>0.
\]

Neither the Sobolev exponent \(p\) nor the effective equidistribution exponent \(\delta\) is made numerically strong enough in the theorem to show \(\eta\geq1/2\). Thus Benoist--Oh supplies a genuine exact-family power saving, but not the requested square-root bound.

\subsection{The sharp hyperbolic frontier}\label{koide-lit-the-sharp-hyperbolic-frontier}

\subsubsection{Good and Chatzakos--Petridis}\label{koide-lit-good-and-chatzakospetridis}

For a cocompact or cofinite Fuchsian group and a fixed hyperbolic conjugacy class \(\mathcal H\), Chatzakos--Petridis define, in equations (1.4)--(1.6),

\[
 M(\mathcal H,X;z)
 =\sum_{1/2<s_j\leq1}
 A(s_j)\widehat u_j u_j(z)X^{s_j},
\]

where \(\lambda_j=s_j(1-s_j)<1/4\) and \(\widehat u_j\) is the period of \(u_j\) along the invariant closed geodesic. Their Theorem 1.1, author manuscript p.~3, restates Good's Theorem 4, p.~116, and proves

\[
 N(\mathcal H,X;z)-M(\mathcal H,X;z)=O(X^{2/3}).
\]

The \(2/3\) arises from the sharp majorant/minorant error

\[
 O(XY^{-1/2}+Y+X^{1/2}\log X)
\]

after choosing \(Y=X^{2/3}\). Theorems 1.2--1.3 are mean-square results, not pointwise improvements. Conjecture 5.7, author manuscript p.~22, explicitly asks for

\[
 O_\varepsilon(X^{1/2+\varepsilon}).
\]

There are two independent failures of direct transfer to \(R(X)\):

\begin{enumerate}
\def\labelenumi{\arabic{enumi}.}

\item
  the Huber cutoff is radial, whereas \(4a+17n\leq X\) plus the chamber is a nonradial sharp polyhedron; norm comparability does not preserve the leading constant or error because the symmetric difference can have volume of order \(X\);
\item
  even in the radial problem, the theorem subtracts every nonzero small-eigenvalue contribution, not only the linear term.
\end{enumerate}

For the exact Koide sector no spectral expansion establishing actual \(X^{s_j}\) terms has yet been produced. They are therefore a necessary obstruction to check, not terms that this audit claims are already known to occur. A proof of \(R(X)-C_0X=O(X^{1/2+\varepsilon})\) would have to show absence or vanishing of all relevant complementary-spectrum coefficients, or explicitly subtract them and then prove they vanish in the target formulation.

\subsubsection{\texorpdfstring{Current sharp refinements remain at \(2/3\)}{Current sharp refinements remain at 2/3}}\label{koide-lit-current-sharp-refinements-remain-at-23}

Voskou, Theorem 3 of arXiv:2407.03134v3, proves \(O(X^{2/3})\) for four prescribed orientation/side sectors in a closed-geodesic-to-closed-geodesic problem, with all small-spectrum period terms included in the main expression. Conjecture 1 on the same page proposes \(O_\varepsilon(X^{1/2+\varepsilon})\), while Theorem 4 is a mean-square result.

Lekkas--Petridis, in the 24 August 2026 version of arXiv:2509.12902, Theorem 1.2, likewise prove

\[
 N(X,l)=\sum_{1/2<s_j\leq1}
 \frac2\pi\gamma_1(s_j)|\widehat u_j|^2X^{s_j}
 +O(X^{2/3})
\]

for a closed-geodesic double-coset count. Their Theorem 1.5 gives

\[
 \frac1X\int_X^{2X}|E(x,l)|^2\,dx\ll X\log^2X,
\]

and Conjecture 1.6 proposes the pointwise \(X^{1/2+\epsilon}\) bound. Their introduction states that Good's \(2/3\) has not been improved for any of the nine point/geodesic/cusp double-coset configurations. This is strong current evidence about the sharp frontier, but neither paper treats the exact Koide height.

\subsubsection{Fixed angular harmonics do not give a sharp interval for free}\label{koide-lit-fixed-angular-harmonics-do-not-give-a-sharp-interval-for-free}

Bourgain--Kontorovich--Sarnak, Theorem 1.1, prove for fixed left and right \(K\)-Fourier modes \(n,k\)

\[
 \sum_{\|\gamma\|<T}e^{2in\theta_1(\gamma)}e^{2ik\theta_2(\gamma)}
 =\mathbf1_{n=0}\mathbf1_{k=0}c_0T^2
  +\sum_jc_j(n,k)T^{2s_j}+O_{n,k}(T^{4/3}).
\]

On the adjoint orbit, affine height is quadratic in the group norm, so \(T^{4/3}\) again corresponds to \(X^{2/3}\). Recovering a sharp angular interval requires summing many Fourier modes, while the theorem's implied constants depend on \(n,k\). It is not a uniform desmoothing theorem for this chamber.

\subsection{Smooth spectral counting}\label{koide-lit-smooth-spectral-counting}

Liu--Sarnak's paper is the closest source with a square-root-plus-gap exponent. Its global arithmetic setup and the hypotheses of its analytic proposition must be separated.

The introduction, author PDF p.~2, fixes an indefinite anisotropic integral ternary form \(f\), a nonzero level \(t\), and assumes \(t\det(f)\) is square-free. Thus the paper's affine-sieve theorem, as globally stated, does not cover \(t\det(f)=-108\). Equations (2.3)--(2.4), author PDF p.~4, fix one arithmetic orbit \(\mathcal O\) and a smooth \(K\)-invariant radial weight \(F_T\) satisfying

\[
 F_T=1\quad(|x|\leq T/c_0),
 \qquad
 F_T=0\quad(|x|\geq c_0T),
\]

with \(c_0>1\) fixed. In Fermi coordinates, Proposition 3.2 uses a smooth even function equal to one for \(|\rho|\leq R-1\), zero for \(|\rho|\geq R+1\), and with uniformly bounded derivatives. Unlike the global sieve statement, Proposition 3.2 is formulated for any variable subgroup \(\Gamma<\mathrm{SL}_2(\mathbb R)\) satisfying its uniform local condition (3.1), the spectral gap (2.7), and containing the fixed primitive cyclic group \(\Gamma_\xi\). It proves

\[
 H_R(z)=\frac{X(R)}{\operatorname{Area}(\Gamma\backslash\mathbb H)}
 +O_\varepsilon(e^{(1/2+\theta+\varepsilon)R}),
 \qquad X(R)\asymp e^R.
\]

Equation (3.60) has \(T=\sinh R\), and equation (3.65) transfers this to the one-sheeted affine orbit. Thus the smooth affine error is

\[
 O_\varepsilon(T^{1/2+\theta+\varepsilon}).
\]

Equation (2.7) assumes \(\lambda_1\geq1/4-\theta^2\). The paper invokes Jacquet--Langlands and Kim--Sarnak to take \(\theta=7/64\), giving

\[
 O_\varepsilon(T^{39/64+\varepsilon})
\]

unconditionally in its congruence-arithmetic setting; \(\theta=0\), and
hence square root, is conditional on the generalized
Selberg/Ramanujan-at-infinity conjecture for the relevant
Jacquet--Langlands transfers.

The numerical input itself is Kim--Sarnak, Appendix 2, Proposition 2, printed p.~176. It applies to an arbitrary cuspidal automorphic representation of \(\mathrm{GL}_2/\mathbb Q\): for a spherical real component, \(|\Re\mu_{j,\infty}|\leq7/64\), independently of finite conductor. Jacquet--Langlands preserves the real component because the quaternion algebra is split at infinity. Chu--Li, proof of Corollary 1.2, equation (5.2), printed p.~5126, records the resulting

\[
 \lambda_1\geq\frac{975}{4096}
\]

for every congruence arithmetic Fuchsian group. Thus the nonmaximal or
non-Eichler structure of the exact finite order does not invalidate
\(s\leq39/64\).  It changes the finite compact-open subgroup and hence which
automorphic representations have fixed vectors and the dimensions of those
fixed-vector spaces; it does not change this archimedean bound. This
validates the spectral-gap input for the exact lattice. It does not provide
sharp-boundary control or force the periods of possible
complementary-spectrum eigenfunctions to vanish.

For the exact order lattice \(\Gamma=\mathcal O^1\), each primitive trace-four conjugacy class supplies the cyclic subgroup in Proposition 3.2, and there are only finitely many such classes. Hence the proposition and equation (3.65) do give an unconditional \(39/64\) estimate after summing the corresponding \textbf{smooth standard-radial} counts. The square-free hypothesis is not an obstruction to this isolated analytic application.

This still does not desmooth to the target. A transition of fixed width in \(R\) is a fixed multiplicative annulus in \(T\), whose invariant volume is of order \(T\), not \(T^{1/2}\). Shrinking the transition makes the derivative/Sobolev constants in Proposition 3.2 grow, and the paper gives no uniform estimate that balances this loss at square-root scale. It also has no sharp coordinate chamber, and its radial height is not \(4a+17n\). These geometric failures remain even after the exact lattice and finite class sum have been supplied.

The same distinction applies to any informal statement that a tempered spectral component has square-root size. Duke--Rudnick--Sarnak's discussion on p.~155 estimates individual tempered frequencies; it does not control the complete spectral sum evaluated at a point with a shrinking approximation to a sharp characteristic function.

\subsection{General effective sector and full-lattice theorems}\label{koide-lit-general-effective-sector-and-full-lattice-theorems}

\subsubsection{Gorodnik--Nevo's explicit generic exponent}\label{koide-lit-gorodniknevos-explicit-generic-exponent}

Gorodnik--Nevo Definition 3.1 defines

\[
 n_e(2)=1,
 \qquad
 n_e(p)=\text{the least even integer at least }p/2\quad(p>2).
\]

Their Section 8 considers

\[
 S_t(\Phi,\Psi)
 =\{g:\log\|gv_0\|\leq t\}\cap\Phi A^+\Psi.
\]

Theorem 8.1, author PDF p.~34, assumes \(L^{p+}\) matrix coefficients and Lipschitz-well-rounded fixed angular sets; Corollary 8.3, p.~35, passes to an orbit when \((\Gamma\cap H)\backslash H\) is compact. The exponent is

\[
 \zeta=\frac{1}{2n_e(p)(1+3\dim G)}.
\]

For \(\dim G=3\),

\[
 \zeta=\frac1{20n_e(p)}.
\]

Even tempered spectrum has \(p=2\), \(n_e(p)=1\), and hence absolute error \(O(X^{19/20+\epsilon})\). For the exact congruence arithmetic lattice, Jacquet--Langlands and Kim--Sarnak give \(s\leq39/64\). The corresponding matrix coefficients are \(L^{p+}\) for \(p=64/25\), so \(n_e(p)=2\) and the theorem gives only \(O(X^{39/40+\epsilon})\).

This numerical calculation is an upper benchmark for the published generic polar-sector theorem, not a claim that Corollary 8.3 literally counts \(R(X)\). A fixed polar sector contains every radial parameter in fixed polar directions. An arbitrary affine linear cone and its finite-radius polyhedral height need an exact polar decomposition or a quantitatively controlled symmetric-difference argument; all integral orbits must then be included. None of those identifications appears in the theorem statement.

Gorodnik--Nevo themselves note after Corollary 5.3, author PDF p.~23, that their sharp-set counting mechanism remains worse than square-root volume error even for tempered spectrum, whereas square-root quality is available there only for a smooth weighted problem.

\subsubsection{Browning--Gorodnik and subsequent full-lattice work}\label{koide-lit-browninggorodnik-and-subsequent-full-lattice-work}

Browning--Gorodnik Theorem 2.1 gives, for a Euclidean ball \(B_H\) in one congruence orbit,

\[
 |\Gamma_\ell y\cap B_H|
 =\frac{m_{Z_{y,\ell}}(Z_{y,\ell})}{m_{X_\ell}(X_\ell)}m_Y(B_H)
 +O\!\left(\ell^{\dim H+\dim G}m_Y(B_H)^{1-\varrho}\right)
\]

for some unspecified \(\varrho>0\). Corollaries 2.5--2.6 sum adelic and integral orbits and retain the Brauer/Kottwitz density, at the cost of \(\ell^{\dim H+2\dim G}\). This is a useful exact-lattice precedent, but its archimedean set is the full norm ball, not the present sector, and no square-root exponent is asserted.

Cao--Huang Theorem 1.5 proves arithmetic purity of the Hardy--Littlewood property for an anisotropic ternary affine quadric with \(-m\det(q)\) nonsquare. Its Theorem 3.4 invokes Nevo--Sarnak, Gorodnik--Nevo, and Browning--Gorodnik to furnish an unspecified positive exponent, denoted \(\beta_X\) there with \(X\) meaning the affine variety; write it here as \(\beta_{\mathscr X}>0\) to distinguish it from the height variable. It is not a sharp-sector remainder theorem.

The current Cao--Huang--Zhang preprint, arXiv:2509.23320, Setting 2.1 and Theorem 2.1, gives a congruence-uniform power saving

\[
 O\!\left(\mathbf N(\mathfrak l)^\kappa
 \operatorname{vol}(\mathbb B_S(N))^{1-\delta}\right)
\]

for standard product Euclidean/max height balls and a compact piecewise-smooth factor. The exponent \(\delta\) is unspecified, and the unbounded real factor in the target is neither compact nor their standard Euclidean ball. This preprint extends the general effective framework but does not improve the conclusion of the direct Benoist--Oh argument in \cref{koide-lit-definition-13.1-for-the-actual-b_x,koide-lit-finite-orbit-passage-and-main-term-normalization}.

\subsection{Results in different arithmetic regimes}\label{koide-lit-results-in-different-arithmetic-regimes}

\subsubsection{Square discriminant and divergent stabilizer}\label{koide-lit-square-discriminant-and-divergent-stabilizer}

Oh--Shah Theorem 1.2 assumes

\[
 \operatorname{vol}((H\cap\Gamma)\backslash H)=\infty
\]

and obtains a smoothed \(cT\log T+O(T)\) count. Its sharp version has only a logarithmic relative error. Theorem 6.1 allows an open cone and arbitrary norm, but retains the logarithmic main-factor characteristic of its divergent-stabilizer regime; that regime is excluded here.

Xu--Zhang, \emph{Quarterly Journal of Mathematics} 74 (2023), 659--685, assumes \(-m\det(f)\) is a square over the number field. Huang, arXiv:2412.03350, Theorem 1.1, assumes the same split condition and a fixed smooth compact weight; its principal error is \(B(\log B)^{1983/1984}\) in a \(B\log B\) count. Zhang, arXiv:2603.21080, Section 1.1 and Theorem 1.1, again assumes \(-m\det(q)\in(k^\times)^2\) and proves a smooth \(K_0\)-invariant expansion with an error \(O(T^{l_1+l_2-\delta})\), \(0<\delta<1/2\); Remark 1.2 obtains an authentic count for \(\Gamma=\mathrm{SL}_2(\mathcal O_k)\). These are important current results, but all lie in the rationally split torus regime excluded by \(108\notin\mathbb Q^{\times2}\), and none uses the target polyhedral height.

\subsubsection{Varying level, isotropic forms, and special rotational quadrics}\label{koide-lit-varying-level-isotropic-forms-and-special-rotational-quadrics}

Huang, arXiv:2412.05903, studies \(F(x)=m_0N\) with \(N=p_0^{2h}\to\infty\), a smooth compactly supported real weight \(w(x/\sqrt N)\), and prescribed local conditions. Theorems 1.1--1.2 have a \(\sqrt N/(\log N)^A\) error. This is a Linnik-type varying-level theorem, not a fixed affine shell.

Li--Liu, arXiv:2501.15033v2 (18 April 2025), Theorems 1.1--1.2, prove new almost-prime conclusions for anisotropic ternary equations. Their global assumption remains \(t\det(f)\) square-free. The quantitative counting input in Lemma 2.3, pp.~5--6, is

\[
 O_\varepsilon(d^{1+\varepsilon}T^{1/2+\theta+\varepsilon})
\]

for the same kind of smooth radial orbit weight inherited from Liu--Sarnak. It is a current application of the \(7/64\) gap, not a new unsmoothed sector theorem.

Lowry-Duda's thesis, arXiv:1704.02376, Theorem 5.5.4, treats the special rotational form \(x_1^2+x_2^2=x_3^2+h\) in a full Euclidean ball. Written in ordinary radius \(T\), its three-dimensional sharp error is \(O(T^{21/22+\epsilon})\). The homogeneous form is \(\mathbb Q\)-isotropic, and neither the lattice nor the height sector is the target one.

Results on values of irrational ternary forms, shrinking sets while the discriminant varies, or positive-definite ternary representation numbers answer different problems. Lindelöf or generalized Lindelöf hypotheses occurring there do not yield the present fixed-level pointwise sector bound.

\subsection{Spectral-gap, GRH, smoothing, and averaging taxonomy}\label{koide-lit-spectral-gap-grh-smoothing-and-averaging-taxonomy}

\begin{longtable}[]{@{}
  >{\raggedright\arraybackslash}p{0.31\linewidth}
  >{\raggedright\arraybackslash}p{0.31\linewidth}
  >{\raggedright\arraybackslash}p{0.31\linewidth}@{}}
\toprule\noalign{}
\begin{minipage}[b]{\linewidth}\raggedright
Classification
\end{minipage} & \begin{minipage}[b]{\linewidth}\raggedright
Source result
\end{minipage} & \begin{minipage}[b]{\linewidth}\raggedright
Why it does not prove the target
\end{minipage} \\
\midrule\noalign{}
\endhead
\bottomrule\noalign{}
\endlastfoot
Unconditional, exact sharp family & Benoist--Oh Theorem 13.2 plus the direct Definition 13.1 verification & Gives only \(O(X^{1-\eta})\) with \(\eta>0\) unspecified \\
Unconditional, exact norm/cone, full lattice & Gorodnik--Oh--Shah Theorem 1.6, with the exact hypotheses proved in \cref{koide-count-lem-genuine-norm,koide-count-lem-gos-hypotheses} & Qualitative \(C_0X+o(X)\) only \\
Unconditional, sharp radial closed-geodesic count & Good / Chatzakos--Petridis & \(O(X^{2/3})\), different height, and main expression includes small-spectrum terms \\
Unconditional, smooth radial count & Liu--Sarnak Proposition 3.2 with \(\theta=7/64\), summed over the exact finite conjugacy-class decomposition & \(39/64\) for the exact lattice, but fixed multiplicative smoothing and a different radial height \\
Selberg/Ramanujan conditional & Liu--Sarnak with \(\theta=0\) & Square-root only for the smooth radial problem; still not the exact count \\
GRH/Lindelöf conditional & None among the source versions in \cref{koide-lit-primary-references-and-exact-locators} proves this fixed sharp-sector estimate & The occurrences listed there concern varying levels/forms, shrinking-set equidistribution, or lower/Omega results \\
Smoothed & Liu--Sarnak; Oh--Shah; Huang I; Zhang 2026 & Smoothing or split arithmetic changes the theorem; no square-root desmoothing supplied \\
Averaged/mean square & Chatzakos--Petridis Theorems 1.2--1.3; Voskou Theorem 4; Lekkas--Petridis Theorem 1.5 & Does not imply a uniform pointwise estimate at every \(X\) \\
Omega/lower-bound & Chatzakos (2019), including a cocompact geodesic-average \(\Omega(X^{1/2}\log\log\log X)\) & Evidence for scale, not an upper bound \\
Isotropic or split & Oh--Shah; Xu--Zhang; Huang I; Zhang 2026; Lowry-Duda & Different stabilizer, quotient, main order, form, or level \\
Varying level & Huang II & Counts \(F=m_0p_0^{2h}\), not \(F=-2\) \\
\end{longtable}

In particular, ``assuming GRH'' is not a substitute for the missing theorem. The conditional hypothesis that occurs naturally in the closest smooth spectral result is the Selberg/Ramanujan eigenvalue conjecture, and even it only reaches square root before the sharp-boundary unsmoothing problem is addressed.

\subsection{Citation-ready conclusion}\label{koide-lit-citation-ready-conclusion}

Subject to the bounded-search qualification in
\cref{koide-lit-audit-limitations}, the following paragraph records the
conclusion supported by the enumerated source versions:

\begin{quote}
Among the source versions enumerated in \cref{koide-lit-primary-references-and-exact-locators}, none implies the unconditional sharp estimate \(R(X)-C_0X=O_\varepsilon(X^{1/2+\varepsilon})\) for the fixed ternary quadric \(a^2-3q^2-18n^2=-2\) in the chamber \(a>0\), \(q,n\geq0\), \(a+4n>q\), with height \(4a+17n\leq X\). Gorodnik--Oh--Shah, Theorem 1.6, applies to the exact admissible generic cone and an arbitrary norm---with those exact hypotheses proved in \cref{koide-count-lem-genuine-norm,koide-count-lem-gos-hypotheses}---but gives only a qualitative asymptotic. Benoist--Oh, Definition 13.1 and Theorem 13.2, applies after the direct effective-well-roundedness verification in \cref{koide-lit-definition-13.1-for-the-actual-b_x} and, after summing the finitely many integral spin-group orbits, yields \(R(X)=C_0X+O(X^{1-\eta})\) for some unspecified \(\eta>0\), not \(\eta=1/2\). In the closest sharp closed-geodesic problem, Good's theorem, as restated and reproved by Chatzakos--Petridis, Theorem 1.1, gives \(O(X^{2/3})\) only after all small-eigenvalue terms \(X^{s_j}\) are included; Chatzakos--Petridis, Conjecture 5.7, states the \(O_\varepsilon(X^{1/2+\varepsilon})\) improvement as a conjecture. Liu--Sarnak, Proposition 3.2, applies class by class to the exact congruence quaternionic lattice and proves the smooth radial bound \(O_\varepsilon(X^{1/2+\theta+\varepsilon})\), unconditionally with \(\theta=7/64\) and conditionally with \(\theta=0\), but its fixed-width smoothing and radial height do not match the exact sharp count. Thus complementary-spectrum terms and sharp nonradial boundary control remain separate unresolved inputs.
\end{quote}

\subsection{Primary references and exact locators}\label{koide-lit-primary-references-and-exact-locators}

\begin{enumerate}
\def\labelenumi{\arabic{enumi}.}

\item
  Yves Benoist and Hee Oh, ``Effective equidistribution of \(S\)-integral points on symmetric varieties,'' \emph{Annales de l'Institut Fourier} \textbf{62} (2012), 1889--1942. \href{https://doi.org/10.5802/aif.2738}{DOI}, \href{https://www.numdam.org/item/AIF_2012__62_5_1889_0/}{official article/PDF}, \href{https://arxiv.org/abs/0706.1621}{arXiv}. Standing setting pp.~1892--1893; finite orbit statement p.~1892; Theorem 1.3 p.~1894; Theorem 12.5 p.~1929; Definition 13.1 p.~1931; Theorem 13.2 pp.~1931--1932; Proposition 14.2 p.~1933; ternary Example (A) p.~1896.
\item
  Alexander Gorodnik, Hee Oh, and Nimish Shah, ``Integral points on symmetric varieties and Satake compactifications,'' \emph{American Journal of Mathematics} \textbf{131} (2009), 1--57. \href{https://doi.org/10.1353/ajm.0.0034}{DOI}, \href{https://arxiv.org/abs/math/0610497}{arXiv}, \href{https://gauss.math.yale.edu/~ho2/doc/satake_ajm.pdf}{author PDF}. Admissible/generic definitions and Theorem 1.6 pp.~4--5; Theorem 1.18 pp.~8--9.
\item
  Alexander Gorodnik, Hee Oh, and Nimish Shah, ``Strong wavefront lemma and counting lattice points in sectors,'' \emph{Israel Journal of Mathematics} \textbf{176} (2010), 419--444. \href{https://doi.org/10.1007/s11856-010-0035-8}{DOI}, \href{https://arxiv.org/abs/0708.3188}{arXiv}. Theorem 1.13, author PDF pp.~4--5.
\item
  Alexander Gorodnik and Amos Nevo, ``Counting lattice points,'' \emph{Journal für die reine und angewandte Mathematik} \textbf{663} (2012), 127--176. \href{https://doi.org/10.1515/CRELLE.2011.096}{DOI}, \href{https://arxiv.org/abs/0903.1515}{arXiv}, \href{https://www.math.uzh.ch/gorodnik/papers/GN2-final.pdf}{author PDF}. Theorem 1.9 pp.~4--5; Definition 3.1 p.~13; Theorem 4.5 pp.~17--18; Theorem 8.1 and Corollary 8.3 pp.~34--35; Proposition 8.4 p.~35.
\item
  William Duke, Zeév Rudnick, and Peter Sarnak, ``Density of integer points on affine homogeneous varieties,'' \emph{Duke Mathematical Journal} \textbf{71} (1993), 143--179. \href{https://doi.org/10.1215/S0012-7094-93-07107-4}{DOI}, \href{https://www.math.ucla.edu/~wdduke/preprints/density.pdf}{author PDF}. Theorem 1.2 p.~145; Example 1.5 pp.~146--147; discussion of individual tempered frequencies p.~155; Theorem 3.5 p.~161.
\item
  Alex Eskin and Curtis McMullen, ``Mixing, counting, and equidistribution in Lie groups,'' \emph{Duke Mathematical Journal} \textbf{71} (1993), 181--209. \href{https://doi.org/10.1215/S0012-7094-93-07108-6}{DOI}, \href{https://people.math.harvard.edu/~ctm/papers/home/text/papers/mixing/mixing.pdf}{author PDF}. Theorem 1.4, author PDF p.~3; Theorems 6.1 and 6.4 pp.~30--31; Open Question 1 p.~35.
\item
  Jianya Liu and Peter Sarnak, ``Integral points on quadrics in three variables whose coordinates have few prime factors,'' \emph{Israel Journal of Mathematics} \textbf{178} (2010), 393--426. \href{https://doi.org/10.1007/s11856-010-0069-y}{DOI}, \href{https://web.math.princeton.edu/facultypapers/sarnak/few-final.pdf}{author PDF}. Arithmetic hypothesis p.~2; equations (2.3)--(2.8) pp.~4--5; Proposition 3.2 pp.~10--11; equations (3.60) and (3.65) pp.~15--16.
\item
  Anton Good, \emph{Local Analysis of Selberg's Trace Formula}, Lecture Notes in Mathematics 1040, Springer, 1983. \href{https://doi.org/10.1007/BFb0073074}{DOI}. Theorem 4, p.~116. The theorem statement and a new proof were checked through Chatzakos--Petridis (2016) because the monograph text was not directly available in this audit.
\item
  Dimitrios Chatzakos and Yiannis N. Petridis, ``The hyperbolic lattice point problem in conjugacy classes,'' \emph{Forum Mathematicum} \textbf{28} (2016), 981--1003. \href{https://doi.org/10.1515/forum-2015-0102}{DOI}, \href{https://arxiv.org/abs/1504.01307}{arXiv}, \href{https://discovery.ucl.ac.uk/1477487/7/Petridis_forum-2015-0102.pdf}{author manuscript}. Equations (1.4)--(1.6) and Theorem 1.1 p.~3; Theorems 1.2--1.3 pp.~3--4; Conjecture 5.7 p.~22.
\item
  Dimitrios Chatzakos, ``Mean value results and \(\Omega\)-results for the hyperbolic lattice point problem in conjugacy classes,'' \emph{Revista Matemática Iberoamericana} \textbf{35} (2019), 1123--1152. \href{https://doi.org/10.4171/RMI/1080}{DOI}, \href{https://arxiv.org/abs/1610.01462}{arXiv}. Mean-value and Omega theorems, not a pointwise upper bound.
\item
  Jean Bourgain, Alex Kontorovich, and Peter Sarnak, ``Sector estimates for hyperbolic isometries,'' \emph{Geometric and Functional Analysis} \textbf{20} (2010), 1175--1200. \href{https://doi.org/10.1007/s00039-010-0092-5}{DOI}, \href{https://arxiv.org/abs/1001.4541}{arXiv}. Theorem 1.1, p.~2.
\item
  Hee Oh and Nimish Shah, ``Limits of translates of divergent geodesics and integral points on one-sheeted hyperboloids,'' \emph{Israel Journal of Mathematics} \textbf{199} (2014), 915--931. \href{https://doi.org/10.1007/s11856-013-0063-2}{DOI}, \href{https://arxiv.org/abs/1104.4988}{arXiv}. Theorem 1.2 p.~3; Corollary 1.4 pp.~3--4; Theorem 6.1 pp.~13--14.
\item
  Tim D. Browning and Alexander Gorodnik, ``Power-free values of polynomials on symmetric varieties,'' \emph{Proceedings of the London Mathematical Society} \textbf{114} (2017), 1044--1080. \href{https://doi.org/10.1112/plms.12030}{DOI}, \href{https://arxiv.org/abs/1606.06342}{arXiv}. Theorem 2.1 p.~11; Corollaries 2.5--2.6 pp.~24--25.
\item
  Yang Cao and Zhizhong Huang, ``Arithmetic purity of the Hardy--Littlewood property and geometric sieve for affine quadrics,'' \emph{Advances in Mathematics} \textbf{398} (2022), 108236. \href{https://doi.org/10.1016/j.aim.2022.108236}{DOI}, \href{https://arxiv.org/abs/2003.07287}{arXiv}. Theorem 1.5 p.~4; Proposition 2.6 pp.~12--13; Theorem 3.4 pp.~21--22.
\item
  Yang Cao, Zhizhong Huang, and Runlin Zhang, ``Effective equidistribution, arithmetic purity of strong approximation, and geometric sieve for affine quadrics,'' \href{https://arxiv.org/abs/2509.23320}{arXiv:2509.23320}. Setting 2.1 and Theorem 2.1; arXiv HTML retrieved 4 September 2026.
\item
  Mikhail Borovoi, ``On representations of integers by indefinite ternary quadratic forms,'' \emph{Journal of Number Theory} \textbf{90} (2001), 281--293. \href{https://doi.org/10.1006/jnth.2001.2662}{DOI}, \href{https://arxiv.org/abs/math/0006141}{arXiv}, \href{https://www.math.tau.ac.il/~borovoi/papers/reprfin.pdf}{author PDF}. Theorem 1.4.1 and Proposition 1.5.1; qualitative local/Brauer--Manin coefficient, no error exponent.
\item
  Fei Xu and Runlin Zhang, ``Counting integral points on indefinite ternary quadratic equations over number fields,'' \emph{Quarterly Journal of Mathematics} \textbf{74} (2023), 659--685. \href{https://doi.org/10.1093/qmath/haac039}{DOI}, \href{https://arxiv.org/abs/2103.10707}{arXiv}. Square-discriminant hypothesis in the main theorem.
\item
  Zhizhong Huang, ``Quantitative strong approximation for ternary quadratic
  forms I.''

  \href{https://arxiv.org/abs/2412.03350}{arXiv:2412.03350}; Theorem 1.1;
  square-discriminant, smooth fixed-weight \(B\log B\) count.
\item
  Zhizhong Huang, ``Quantitative strong approximation for ternary quadratic
  forms II.''

  \href{https://arxiv.org/abs/2412.05903}{arXiv:2412.05903}; Theorems 1.1--1.2;
  varying level \(m_0p_0^{2h}\) and smooth compact weight.
\item
  Runlin Zhang, ``Effective count of integer points on ternary affine quadrics and effective equidistribution,'' \href{https://arxiv.org/abs/2603.21080}{arXiv:2603.21080}. Section 1.1 and Theorem 1.1; split-stabilizer smooth count; Remark 1.2 on the authentic radial count.
\item
  Marios Voskou, ``Refined Counting of Geodesic Segments in the Hyperbolic
  Plane.''

  \href{https://arxiv.org/abs/2407.03134}{arXiv:2407.03134v3}; Theorems 3--4
  and Conjecture 1, pp.~3--4.
\item
  Dimitrios Lekkas and Yiannis N. Petridis, ``A relative trace formula and counting geodesic arcs in the hyperbolic plane,'' \href{https://arxiv.org/abs/2509.12902}{arXiv:2509.12902}, version dated 24 August 2026. Theorem 1.2, Theorem 1.5, and Conjecture 1.6.
\item
  Henry H. Kim, ``Functoriality for the exterior square of \(\mathrm{GL}_4\) and the symmetric fourth of \(\mathrm{GL}_2\),'' with Appendix 2 by Henry H. Kim and Peter Sarnak, ``Refined estimates towards the Ramanujan and Selberg conjectures,'' \emph{Journal of the American Mathematical Society} \textbf{16} (2003), 139--183. \href{https://doi.org/10.1090/S0894-0347-02-00410-1}{DOI}, \href{https://www.ams.org/journals/jams/2003-16-01/S0894-0347-02-00410-1/S0894-0347-02-00410-1.pdf}{official AMS PDF}. Appendix 2, Proposition 2, printed p.~176: \(|\Re\mu_{j,\infty}|\leq7/64\) for a spherical real component of any cuspidal \(\mathrm{GL}_2/\mathbb Q\) representation, and the displayed consequence \(\lambda_1\geq975/4096\).
\item
  Michelle Chu and Han Li, ``Small generators of cocompact arithmetic Fuchsian groups,'' \emph{Proceedings of the American Mathematical Society} \textbf{144} (2016), 5121--5127. \href{https://doi.org/10.1090/proc/13177}{DOI}, \href{https://arxiv.org/abs/1603.00926}{arXiv}. Section 3.1, p.~5122, defines congruence arithmetic Fuchsian groups; proof of Corollary 1.2, equation (5.2), p.~5126, records \(\lambda_1\geq975/4096\) for every such group via Jacquet--Langlands and Kim--Sarnak.
\item
  Hervé Jacquet and Robert P. Langlands, \emph{Automorphic Forms on \(\mathrm{GL}(2)\)}, Lecture Notes in Mathematics 114, Springer, 1970. \href{https://doi.org/10.1007/BFb0058988}{DOI}, \href{https://publications.ias.edu/sites/default/files/automorphic-forms-on-gl2_rpl_9.pdf}{IAS-hosted text}. Chapter III, especially §16, supplies the global quaternionic transfer; at split places, including the real place here, the local component is preserved.
\item
  John Voight, \emph{Quaternion Algebras}, Graduate Texts in Mathematics 288, Springer, 2021. \href{https://doi.org/10.1007/978-3-030-56694-4}{DOI}, \href{https://jvoight.github.io/quat-book.pdf}{open author version}. Proposition 15.5.2 (maximal superorders); §23.1.3 (Eichler orders at ramified primes); equations (27.6.6)--(27.6.7) (compact-open finite adelic order/unit group); Theorem 28.2.10 (strong approximation for indefinite \(B^1\)); Definition 38.7.2 (principal congruence groups).
\item
  Jiamin Li and Jianya Liu, ``Anisotropic quadratic equations in three
  variables.''

  \href{https://arxiv.org/abs/2501.15033}{arXiv:2501.15033v2};
  Theorems 1.1--1.2, pp.~2--3; Lemma 2.3, pp.~5--6. The global form-level
  hypothesis is \(t\det(f)\) square-free; the quantitative input is a smooth
  radial, congruence-orbit estimate with
  \[
   O_\varepsilon\!\left(
     d^{1+\varepsilon}T^{1/2+\theta+\varepsilon}
   \right).
  \]
\item
  David Lowry-Duda, \emph{On Some Variants of the Gauss Circle Problem}, \href{https://arxiv.org/abs/1704.02376}{arXiv:1704.02376}. Theorem 5.5.4, thesis pp.~97--98.
\end{enumerate}

\subsection{Audit limitations}\label{koide-lit-audit-limitations}

The claim-by-claim source versions and locators used in this bounded audit
are enumerated in
\cref{koide-lit-primary-references-and-exact-locators}.  The canonical tree
does not preserve byte-identical snapshots or retrieval hashes for those
external publications, so this file does not assert a reproducible local
source corpus.  Good's 1983 book was the sole exception to direct full-text
access: its theorem locator and statement were checked through the explicit
restatement and independent proof in Chatzakos--Petridis (2016) and the
version-dated Lekkas--Petridis treatment.

The statement that no theorem was located is a result about the source
versions enumerated here, not a proof that no unpublished or differently
indexed result exists.  It does establish that none of the enumerated works
may be cited as a proof of the exact unconditional square-root discrepancy.
